\clearpage
\SourceLeaf{227}%
\begin{center}\Large\textsc{[Johann Bolyais Appendix.]}\end{center}
\bigskip
\begin{center}\textsc{Gauss an Gerling. Göttingen, 14. Februar 1832.}\end{center}
\mbox{.\ .\ .\ .\ .\ .\ .} Noch bemerke ich, dass ich dieser Tage eine kleine Schrift aus \textls{Ungarn} über die Nicht-Euklidische Geometrie erhalten habe, worin ich alle \textls{meine eigenen Ideen und Resultate} wiederfinde, mit grosser Eleganz entwickelt, obwohl in einer für jemand, dem die Sache fremd ist, wegen der Concentrirung etwas schwer zu folgenden Form. Der Verfasser ist ein sehr junger österreichischer Officier, Sohn eines Jugendfreundes von mir, mit dem ich 1798 mich oft über die Sache unterhalten hatte, wiewohl damals meine Ideen noch viel weiter von der Ausbildung und Reife entfernt waren, die sie durch das eigene Nachdenken dieses jungen Mannes erhalten haben. Ich halte diesen jungen Geometer v.~\textsc{Bolyai} für ein Genie erster Grösse. \mbox{.\ .\ .\ .\ .\ .}

\bigskip
\begin{center}\rule{.25\textwidth}{.3pt}\end{center}
\bigskip
\begin{center}\textsc{Gauss an Wolfgang von Bolyai. Göttingen, 6. März 1832.}\end{center}
\mbox{.\ .\ .\ .\ .\ .\ .} Jetzt Einiges über die Arbeit Deines Sohnes.

Wenn ich damit anfange, »\textls{dass ich solche nicht loben darf}«: so wirst Du wohl einen Augenblick stutzen: aber ich kann nicht anders; sie loben hiesse mich selbst loben: denn der ganze Inhalt der Schrift, der Weg, den Dein
